
\begin{table}[H]
\centering
\caption{Composición elemental estándar de un microorganismo}
\label{tab:comp_est_micro}
\begin{threeparttable}
{\setstretch{1.25}
\begin{tabular}{cc}\hline
    
    \bfseries{Elemento} & \bfseries{Porcentaje} \\ \hline
    
    C & 46.5\% \\
    H & 6.49\% \\
    O & 31\% \\
    N & 10.85\% \\
    Sales & $\approx$5\% \\
    
    \hline
\end{tabular}
}
\begin{tablenotes}
    \item[] 
\end{tablenotes}
\end{threeparttable}
\end{table}

\section{Ejercicio de balanceo de ecuaciones}

Podemos intentar balancear la siguiente ecuación, la cual representa la producción de ácido acético a partir de glucosa:

\begin{equation*}
    C_{6}H_{12}O_{6} + O_{2} + NH_{3} \Longrightarrow C_{0.039}H_{0.072}O_{0.021}N_{0.069} + C_{6}H_{8}O_{7} + CO_{2} + H_{2}O
\end{equation*}

A primera vista se puede notar que intentar resolver este ejercicio ``al tanteo'' no es realista, por lo que es mejor recurrir al método algebráico, cuyo proceso se realiza a continuación:

\begin{enumerate}
    
    \item Definimos los coeficientes de cada compuesto en base a las letras del abecedario
        \begin{equation*}
            A(C_{6}H_{12}O_{6}) + B(O_{2}) + C(NH_{3}) \Longrightarrow D(C_{0.039}H_{0.072}O_{0.021}N_{0.069}) + E(C_{6}H_{8}O_{7}) + F(CO_{2}) + G(H_{2}O)
        \end{equation*}

    \item Luego, revisamos cuantas veces aparece cada elemento en cada compuesto de la ecuación química, y ordenamos como si fuera una ecuación matemática, juntando las veces que aparece el elemento con la letra del coeficiente.
        \begin{enumerate}
	        \item Por ejemplo, iniciemos con el carbono (C). Tomando que la glucosa ($C_{6}H_{12}O_{6}$) tiene 6 carbonos, y su coeficiente es A, tenemos que el primer término de la ecuación es $\text{6A}$, y repetimos con los demás compuestos.
	        \item Repetimos esto con el resto de elementos ($H$, $O$ y $N$).
	        \item Otro detalle es cambiar la flecha de la ecuación química ($\Longrightarrow$) por un signo de igual ($=$).
        \end{enumerate}

    \item Terminaríamos con las siguientes ecuaciones:
        \begin{align*}
            C &\longrightarrow 6A = 0.39D + 6E + F \\ H &\longrightarrow 12A + 3C = 0.72D + 8E + 2G \\ O &\longrightarrow 6A + 2B = 0.021D + 7E + 2F + G \\ N &\longrightarrow C = 0.0069D
        \end{align*}
\end{enumerate}

Posterior a esto se podrían intentar despejar las diferentes incógnitas como si fueran cualquier otra serie de ecuaciones matemáticas, sin embargo, debemos considerar algo: la cantidad de incógnitas y de ecuaciones.\footnote{Recordar los grados de libertad de la materia Balance de Materia y Energía}

Una regla de las matemáticas es que, para que una serie de ecuaciones sea posible de resolver, estas deben de tener la misma cantidad de incógnitas que de ecuaciones, mientras que si existen más incógnitas que ecuaciones es irresoluble. Si hacemos una equivalencia entre ecuaciones matemáticas y químicas, tenemos que los coeficientes de los compuestos serían las incógnitas, y que cada elemento cuenta con una ecuación. Tomando esto en cuenta, podemos notar que tenemos más incógnitas (7 compuestos/coeficientes) que ecuaciones (elementos). Por ende, este problema es imposible de resolver tal como se presenta.

Afortunadamente, existen valores que pueden ayudar a modificar la ecuación de forma que sea posible de resolver.

\begin{align*}
    Y_{xs} &= \frac{\text{g de biomasa}}{\text{g de sustrato}}
    \\
    \\
    Y_{ps} &= \frac{\text{g de producto}}{\text{g de sustrato}}
    \\
    \\
    x &= \text{Biomasa}
    \\
    s &= \text{Sustrato}
    \\
    p &= \text{Producto}
\end{align*}

El rendimiento refleja la cantidad de ``algo'' obtenido tras la inversión de otro ``algo''. En bioprocesos el elemento invertido es el sustrato ($s$), y se puede obtener biomasa ($x$) y producto ($p$).
Estos valores, obtenidos a través de pruebas de laboratorios, nos pueden ayudar a despejar los coeficientes de la biomasa y del producto (Ácido acético, $C_{6}H_{8}O_{7}$), dejando la ecuación con 4 incógnitas y 4 ecuaciones, o sea, posible de resolver.

\begin{align*}
    \text{Rendimientos} \longrightarrow
    &Y_{xs} = 0.153
    \\
    &Y_{ps} = 0.85
    \\
    \\
    \text{Pesos molares} \longrightarrow
    &PM_{s} = \frac{180g}{mol}
    \\
    \\
    &PM_{x} = \frac{96.5g}{mol}
    \\
    \\
    &PM_{p} = \frac{192g}{mol}
    \\
    \\
    \text{Coeficientes} \longrightarrow
    &\frac{0.153g_{x}}{1g_{s}} \cdot \frac{1mol_{x}}{96g_{x}} \cdot \frac{180g_{s}}{1mol_{s}} \cdot = 0.287\frac{mol_{x}}{mol_{s}}
    \\
    \\
    &\frac{0.85g_{p}}{1g_{s}} \cdot \frac{1mol_{p}}{192g_{p}} \cdot \frac{180g_{s}}{1mol_{s}} \cdot = 0.797\frac{mol_{p}}{mol_{s}}
\end{align*}

Teniendo ahora los coeficientes de la biomasa y del producto, podemos reemplazar en la ecuación química y volver a realizar el balance algebráico

\begin{align*}
    1(C_{6}H_{12}O_{6}) + B(O_{2}) + C(NH_{3}) &\Longrightarrow 0.287(C_{0.039}H_{0.072}O_{0.021}N_{0.069}) + 0.797(C_{6}H_{8}O_{7}) + F(CO_{2}) + G(H_{2}O)
    \\
    \\
    C &\longrightarrow 6(1) = 0.39(0.287) + 6(0.797) + F
    \\
    H &\longrightarrow 12(1) + 3C = 0.72(0.287) + 8(0.797) + 2G
    \\
    O &\longrightarrow 6(1) + 2B = 0.021(0.287) + 7(0.797) + 2F + G
    \\
    N &\longrightarrow C = 0.0069(0.287)
\end{align*}

Despues de realizar los cálculos deberíamos de tener:
\begin{align*}
    1(C_{6}H_{12}O_{6}) + 2.43(O_{2}) &+ (1.8368 \cdot 10^{-3})(NH_{3})
    \\
    &\Longrightarrow
    \\
    0.287(C_{0.039}H_{0.072}O_{0.021}N_{0.069}) &+ 0.797(C_{6}H_{8}O_{7}) + 1.238(CO_{2}) + 2.8(H_{2}O)
\end{align*}